
\documentclass{article}
\usepackage{colortbl}
\usepackage{makecell}
\usepackage{multirow}
\usepackage{supertabular}

\begin{document}

\newcounter{utterance}

\centering \large Interaction Transcript for game `ifeval', experiment `ifeval', episode 240 with Qwen3.5{-}9B.
\vspace{24pt}

{ \footnotesize  \setcounter{utterance}{1}
\setlength{\tabcolsep}{0pt}
\begin{supertabular}{c@{$\;$}|p{.15\linewidth}@{}p{.15\linewidth}p{.15\linewidth}p{.15\linewidth}p{.15\linewidth}p{.15\linewidth}}
    \# & \multicolumn{2}{c}{Player} && \multicolumn{2}{c}{Game Master} \\
    \hline

    \theutterance \stepcounter{utterance}  
    & & & \multicolumn{4}{p{0.6\linewidth}}{
        \cellcolor[rgb]{0.9,0.9,0.9}{
            \makecell[{{p{\linewidth}}}]{
                \texttt{\tiny{[P1$\langle$GM]}}
                \texttt{Write a casual blog post about how the outer solar system is different from the inner solar system, and what that means for the possibility of life. Wrap your entire response with double quotation marks. Your response should contain 17 or more sentences.} \\
            }
        }
    }
     \\ \\

    \theutterance \stepcounter{utterance}  
    & \multicolumn{4}{p{0.6\linewidth}}{
        \cellcolor[rgb]{0.561,0.737,0.561}{
            \makecell[{{p{\linewidth}}}]{
                \texttt{\tiny{[P1$\rangle$GM]}}
                \texttt{"Hey everyone, today we are going to take a quick trip out past our asteroid belt to see how the outer solar system really stacks up against the inner planets. If you look at Mercury, Venus, Earth, and Mars, you will notice they are mostly made of rock and metal, which makes them feel very solid under your feet. In contrast, the giants like Jupiter and Saturn are huge balls of gas and ice that you could never actually stand on. This big difference in composition is the first major clue that tells us the two regions of our solar system are quite distinct. The inner planets are also much closer to the Sun, which means they get a lot more heat and light than their distant cousins. Because of this proximity, the inner worlds tend to have thinner atmospheres or none at all, while the outer giants are wrapped in thick, swirling layers of clouds. Temperature plays a huge role here too, since it gets freezing cold out there near Neptune and beyond. That extreme cold actually helps preserve things that would otherwise evaporate quickly, like water ice and organic compounds. When we think about life, most scientists agree that it needs liquid water to get started, which is why Earth is so special in our neighborhood. The inner planets struggle to keep liquid water on their surfaces because it either boils away or freezes solid depending on how close they are to the star. However, the outer system offers a different kind of hope because of those massive moons orbiting the gas giants. Places like Europa under Jupiter or Enceladus around Saturn might have entire oceans of liquid water hiding beneath their icy crusts. These subsurface oceans are kept warm by the gravitational tugging of their giant parent planets, creating a cozy environment for potential microbes. It is fascinating to imagine tiny alien creatures swimming in the dark depths of an ocean that has never seen sunlight. While we can't see complex life like humans out there just yet, the conditions for simple, single{-}celled organisms seem much more possible in the outer reaches. The rocky inner planets are great for studying the history of our own solar system, but the icy outer worlds might hold the secrets to life elsewhere. We are only just beginning to send our best robots to explore these mysterious places and look for signs of biology. Every new piece of data we get from these distant moons makes me wonder if we are truly alone in this vast universe. So next time you look up at the night sky, remember that the dark, cold corners might be hiding something amazing. Let's keep our eyes on the outer solar system because the future of astrobiology could be written in ice and gas. Thanks for reading, and I hope you enjoyed this little journey through our cosmic backyard!"} \\
            }
        }
    }
    & & \\ \\

    \theutterance \stepcounter{utterance}  
    & & & \multicolumn{2}{p{0.3\linewidth}}{
        \cellcolor[rgb]{0.9,0.9,0.9}{
            \makecell[{{p{\linewidth}}}]{
                \texttt{\tiny{[GM$|$GM]}}
                {'startend:quotation': {}, 'length_constraints:number_sentences': {'relation': 'at least', 'num_sentences': 17}}
            }
        }
    }
    & & \\ \\

    \theutterance \stepcounter{utterance}  
    & & & \multicolumn{2}{p{0.3\linewidth}}{
        \cellcolor[rgb]{0.9,0.9,0.9}{
            \makecell[{{p{\linewidth}}}]{
                \texttt{\tiny{[GM$|$GM]}}
                \texttt{game\_result = WIN} \\
            }
        }
    }
    & & \\ \\

\end{supertabular}
}

\end{document}
